% BEGIN PROFESSOR FEEDBACK PF08: Method - System overview
\section{System overview}
\label{sec:system_overview}

This chapter describes how the detection approach is put into practice. The
method is implemented as a single program that turns raw Ethereum transaction
records into malicious-address scores and category-level comparisons. The
program is not used interactively. It runs a fixed sequence of steps, always in
the same order, so that the same data and the same settings always lead to the
same result.

The program runs in two phases. A \emph{preparation} phase reads the raw data
once, decides which addresses take part in the study, and fixes how they are
divided into a development set and a test set. An \emph{execution} phase then carries
out the actual detection on top of these prepared data. Keeping the two apart
means that the data and the evaluation setup are settled before any model is
trained, so the reported results cannot be shaped, even by accident, by later
changes to the data or the settings. The execution phase is checkpointed as
well. After an expensive subgraph count has been committed successfully, an
interrupted run verifies and reuses it instead of starting that count again.

Overall the program moves from raw transactions to results in eight steps, shown
in Figure~\ref{fig:method_steps}. The first two belong to the preparation phase,
the remaining six to the execution phase.

\begin{enumerate}
  \item \textbf{Loading the data.} The program reads the transaction records and
        the address labels of the chosen dataset. For the main dataset the
        transaction table is very large, with several hundred million transfers,
        so it is read in small batches instead of being held in memory all at
        once. The datasets are described in Section~\ref{sec:data_acquisition}.

  \item \textbf{Preparing the data.} Not every labelled address can be used. The
        program measures how active each address is and keeps only those with
        enough transactions to form a three-edge candidate pattern. It also sets aside a small
        number of extremely active hub addresses, so that a few exchange-sized
        accounts cannot dominate the analysis. The remaining addresses are split
        into a development set, used for tuning, and a separate test set, used
        only for the final measurement. For each address, the program stores just
        the transactions that involve it and writes these tables to
        disk, so that every later step works on the same data.

  \item \textbf{Building the temporal graph.} The prepared transactions are
        loaded into a temporal graph, in which every address is a node and every
        transfer is a directed edge that carries the time at which it happened.
        Repeated transfers between the same two addresses are kept separately,
        because their number and their order define the candidate patterns.
        This step is described in Section~\ref{sec:graph_construction}.

  \item \textbf{Counting the candidate subgraphs.} The program first checks the
        counter against synthetic validation templates. It then counts how
        often each address participates in each retained local role.
        Because an occurrence only counts when its transfers fall within the
        selected time span, the counting is repeated for a range of time windows,
        from five minutes up to 30 days. This step is described in
        Section~\ref{sec:motif_extraction}. Each completed count is stored as an
        immutable checkpoint before the next one begins.

  \item \textbf{Building the fingerprints.} The local motif-role counts of an address form
        its fingerprint, a short list of numbers that records how often the
        address occupied each retained role. From the same transactions the program
        also computes a set of simpler behavioural features, such as how much
        value the address sent and received, how many partners it had, and how
        its activity was spread over time. This produces three views of every
        address: the simpler features on their own, the motif fingerprint on its
        own, and the two together.

  \item \textbf{Running the detection.} For each of the three views, the program
        trains a classifier to tell malicious and benign addresses apart. It
        uses the development set to choose the time window that works best and
        then measures how well the classifier performs on the untouched test
        set. For RQ3 and RQ4 the same idea is used, except that the classifier
        sorts criminal addresses into their behaviour category rather than making
        a malicious-or-benign decision. The features and models are described in
        Section~\ref{sec:detection_methods}.

  \item \textbf{Analysing the results.} The program then answers the research
        questions. It compares the three views to see whether the motif
        fingerprint improves detection over the simpler features (RQ1) and how
        the time window affects its usefulness for detection (RQ2). It then checks
        whether different categories of malicious behaviour have distinct
        fingerprints (RQ3), and whether selecting a separate window for each
        category improves classification over one shared window (RQ4). Each comparison is
        supported by statistical tests, so that a difference is not mistaken for
        chance, and the whole run is repeated under several orderings of
        transactions that share the same timestamp, to confirm that the outcome
        is stable. The evaluation is described in
        Section~\ref{sec:evaluation_protocol}.

  \item \textbf{Visualising the results.} Finally, the program writes the tables
        and figures used in the Results chapter, including the comparison of the
        three views, the performance curve across the different time windows, and
        for RQ3 and RQ4, the category fingerprints and the comparison between
        category-specific and shared windows.
\end{enumerate}

\begin{figure}[p]
  \centering
  \ResGraphic{1.00}{05-method-steps}
  \caption{The eight steps of the pipeline. The first two make up the preparation
  phase. The remaining six cover graph construction, fingerprint extraction,
  detection, analysis, and reporting in the execution phase.}
  \label{fig:method_steps}
\end{figure}

The first two research questions run on the main binary dataset and the last two
on the category-labelled dataset. The following sections describe the most
important of these steps in more detail, starting with the data.
% END PROFESSOR FEEDBACK PF08

